% Helper macros


%%%%%%%%%%%
% Toggles to enable/disable certain parts
%%%%%%%%%%%
\newtoggle{debug}
\newtoggle{title}
\newtoggle{toc}
\newtoggle{notation}
\newtoggle{index}
\ifdefined\Debug
    \toggletrue{debug}
\else
    \togglefalse{debug}
\fi
\ifdefined\noTitle
    \togglefalse{title}
\else
    \toggletrue{title}
\fi
\ifdefined\noTOC
    \togglefalse{toc}
\else
    \toggletrue{toc}
\fi
\ifdefined\noNotation
    \togglefalse{notation}
\else
    \toggletrue{notation}
\fi
\ifdefined\noIndex
    \togglefalse{index}
\else
    \toggletrue{index}
\fi


%%%%%%%%%%%
% Comments and todos
%%%%%%%%%%%
\newcommand{\colorpar}[3]{\colorbox{#1}{\parbox{#2}{#3}}}
% Remark inline filling complete line
\newcommand{\lineremark}[3]{\colorpar{#2}{\linewidth}{\color{#1}#3}\newline}
% Remark in page margin
\newcommand{\marginremark}[3]{\marginpar{\colorpar{#2}{\linewidth}{\color{#1}#3}}}
% Boxes for writing todos, general todos and questions
\newcommand{\towrite}[1]{\lineremark{colorTodo}{colorTodoFill}{\textbf{[WRITE]}: #1}}
\newcommand{\todo}[1]{\colorbox{colorTodoFill}{\textcolor{colorTodo}{\textbf{[TODO]}: #1}}}
\newcommand{\question}[1]{\marginremark{colorQuestion}{colorQuestionFill}{\small{\textbf{[Q?]}: #1}}}
\newcommand{\feedback}[1]{\marginremark{colorFeedback}{colorFeedbackFill}{\small{\textbf{[Fdb]}: #1}}}


%%%%%%%%%%%
% Index
%%%%%%%%%%%
% Create text in margin
\newcommand{\keywordmarginnote}[1]{%
    \marginnote{\setstretch{.6}\parbox[t]{.8\marginparwidth}{\textcolor{colorKeywordMargin}{\textsf{\scriptsize{#1}}}}}%
}%

% Add index entry with optional parent entry
\newcommand{\indexentry}[2][]{%
    \ifthenelse{\equal{#1}{}}{\index{#2}}{\index{#1!#2}}%
}%
\newcommand{\indexentrybf}[2][]{\indexentry[#1]{#2|textbf}}
%
% Add text with index entry
% Usage:
% \indextext[parent]{text}
% \indextextpl[parent]{text}{indextext}
\newcommand{\indextextpl}[3][]{\emph{#2}\indexentry[#1]{#3}}
\newcommand{\indextext}[2][]{\indextextpl[#1]{#2}{#2}}

% Create keyword with margin text and index entry
% Usage:
% \keyword[parent]{text/margintext/indextext}
% \keywordpl[parent]{text}{margintext/indextext}
% \keywordindex[parent]{text}{margintext}{indextext}
% - parent (if provided) is the parent index
% - text is the name used in the text
% - margintext is the name of the keyword used in the margin
% - indextext is the name used in the index
\newcommand{\keywordindex}[4][]{%
    \indextextpl[#1]{\textcolor{colorKeywordText}{#2}}{#4|textbf}% create text and index
    \ifthenelse{\equal{#1}{}}{\keywordmarginnote{#3}}{\keywordmarginnote{$\cdot$#3}}% create margin
}%
\newcommand{\keywordpl}[3][]{\keywordindex[#1]{#2}{#3}{#3}}
\newcommand{\keyword}[2][]{\keywordpl[#1]{#2}{#2}}

\iftoggle{index}{
    % Index
    % Style file is added in .latexmkrc
    \makeindex[title={Index\label{index}}]
    \indexsetup{level=\chapter,toclevel=chapter}
}{}

%%%%%%%%%%%
% Glossary (for list of symbols)
%%%%%%%%%%%

% Create topic for list of symbols
\newcommand{\newsymboltopic}[2]{%
    \newglossaryentry{#1}{%
        name={#2},
        description={},
        type=symbols,
        parent={}
    }%
    %\gls{#1}%
}

% Create symbol entry for list of symbols.
% Usage:
% \newglosssymbol{parent}{label}{symbol}{description}
% The 'parent' label must reference a topic created with \newsymboltopic
\newcommand{\newglosssymbol}[4]{%
    \newglossaryentry{#1:#2}{%
        name=\ensuremath{#3}, % symbol needs to be in math mode
        description={#4},
        type=symbols,
        %symbol={#3},
        parent={#1}
    }%
}
% Display symbol at current position.
% Usage:
% \glosssymbol{parent}{label}
\newcommand{\glosssymbol}[2]{%
    \gls{#1:#2}%
}

% Create symbol for list of symbols and display it at current position.
% Usage:
% \newsymbol{parent}{label}{symbol}{description}
% The 'parent' label must reference a topic created with \newsymboltopic
\newcommand{\newsymbol}[4]{%
    \newglosssymbol{#1}{#2}{#3}{#4}%
    \glosssymbol{#1}{#2}%
}

% Define a new glossary type
\iftoggle{notation}{
    % Glossary
    \makeglossaries
}{}



%%%%%%%%%%%
% Tikz
%%%%%%%%%%%

% External tikz
% must be loaded after \makeindex, otherwise it will not compile
% remove command \includepdf which otherwise leads to compilation issues
\tikzexternalize[prefix=tikz_extern/,optimize command away=\includepdf]

% Overwrite compile command to allow setting of output directory
\makeatletter
    \tikzset{%
        external/shell escape={-shell-escape\space-output-directory="\OutputDir/"},%
        /pgf/images/include external/.code={%
            \includegraphics{\OutputDir/#1}%
        },%
    }
\makeatother

% Tikz picture (without float environment)
% Arguments:
% 1. Name which will be used for filename
% 2. Name which will be used for label
% 3. Input filename in tikz/ directory
% 4. Caption
\newcommand{\tikzPicture}[4]{%
    \tikzsetnextfilename{#1}
    \input{figures/#3}
    \caption{#4}
    \label{#2}%
}
% Tikz figure
% Arguments:
% 1. Name which will be used for filename and label
% 2. Input filename in tikz/ directory
% 3. Caption
% 4. Positioning
\newcommand{\tikzFigurePos}[4]{%
    \begin{figure}[#4]%
        \centering
        \tikzPicture{#1}{fig:#1}{#2}{#3}
    \end{figure}%
}%
\newcommand{\tikzFigure}[3]{%
    \tikzFigurePos{#1}{#2}{#3}{t}%
}%
% Tikz subfigure
% Arguments:
% 1. Name which will be used for filename and label
% 2. Input filename in tikz/ directory
% 3. Caption
% 4. Width
\newcommand{\tikzSubfigure}[4]{%
    \begin{subfigure}[b]{#4}
        \centering
        \tikzPicture{#1}{subfig:#1}{#2}{#3}
    \end{subfigure}%
}


%%%%%%%%%%%
% Table
%%%%%%%%%%%

% Table
% Arguments:
% 1. Name which will be used for label
% 2. Input filename in tables/ directory
% 3. Caption
\newcommand{\newTable}[3]{%
    \begin{table}[t]%
        \centering
        \caption{#3}
        \label{tab:#1}
        \input{tables/#2}%
    \end{table}%
}%
% Table with notes
\newcommand{\newTableNotes}[3]{%
    \begin{table}[t]%
        \centering
        \begin{threeparttable}[t]%
            \centering
            \caption{#3}
            \label{tab:#1}
            \input{tables/#2}%
        \end{threeparttable}%
    \end{table}%
}%


%%%%%%%%%%%
% Algorithm
%%%%%%%%%%%

% Algorithm
% Arguments:
% 1. Name which will be used for label
% 2. Input filename in algorithms/ directory
% 3. Caption
\newcommand{\newAlgorithm}[3]{%
    \begin{algorithm}[t]%
        \centering
        \caption{#3}
        \label{alg:#1}
        \input{algorithms/#2}%
    \end{algorithm}%
}%


%%%%%%%%%%%
% Misc
%%%%%%%%%%%

% Highlight text
\newcommand{\highlight}[1]{\textcolor{colorHighlight}{#1}}

% Chapter thumbs
\newcommand{\partThumbs}[1]{\chapterthumbsoff\part{#1}\chapterthumbson}

% Full reference including name
\newcommand*{\cfullref}[1]{\hyperref[{#1}]{\Cref*{#1} (\nameref{#1})}}
\newcommand*{\vfullref}[1]{\hyperref[{#1}]{\Vref*{#1} (\nameref{#1})}}

% Store text to use later
% Arguments:
% 1. Label to reference the stored text
% 2. Text to store
\newcommand*{\rememberText}[2]{\Copy{rem:#1}{#2}}
% Restate remembered text
% Argument:
% - Label used to reference the stored text
\newcommand*{\restateText}[1]{\Paste{rem:#1}}

% Restate theorem/lemma/... in appendix
% Arguments:
% 1. Label used to reference the theorem/lemma/...
% 2. Label used to reference the stored text of the theorem/lemma/... The proof will get the same label with prefix "proof:"
\newcommand*{\appendixRestate}[2]{\subsection{Proof of~\cref{#1}}\label{proof:#2}
We have to show~\vref{#1}:
\begin{restatebox}%
    \restateText{#2}%
\end{restatebox}%
}
